% chapters/05-status-security.tex -- part B. 상태 코드와 보안.
% Macros come from part A's preamble.tex (\term \eng \code \decision \commit \cfs \todo inbrief callout).
% \code is verbatim-like: not inside \caption, \section, node text. Use \texttt there.
\chapter{상태 코드와 보안}\label{ch:status-security}

% src: deploy/traffic/review/03-status-codes.md §목표, deploy/traffic/review/02-blocking.md §K §N
\paragraph{결과.}
응답 코드가 실제 의미를 말하게 했다.
거절은 4xx로 답하고, 서버 고장만 5xx로 답한다.
공개 쓰기 경로 하나는 비밀 헤더 없이는 DB에 닿지 못한다.
\Cref{tab:05-summary}가 이 장의 변경 여섯 개를 한 줄씩 보인다.

\begin{table}[H]
  \caption{이 장의 변경 요약. 전과 후, 반영한 커밋(commit).}\label{tab:05-summary}
  % src: deploy/traffic/review/03-status-codes.md §G §L §제안; deploy/traffic/review/02-blocking.md §N §K
  \small
  \begin{tabularx}{\linewidth}{l Y Y l}
    \toprule
    항목 & 전 & 후 & 커밋 \\
    \midrule
    G. CORS 거절 & 500, ERROR 2줄, DB INSERT 1건 & 403, WARN 1줄, DB 쓰기 없음 & \commit{4d137f2} \\
    \midrule
    L. error-logs 쓰기 & 누구나 \code{curl}로 INSERT 가능 & \code{x-origin-key}가 맞을 때만 통과, 아니면 403 & \commit{38b4c12} \\
    \midrule
    제안 1. 서버 종류 헤더 & \code{x-powered-by}가 보임 & 헤더 없음 & \commit{38b4c12} \\
    \midrule
    제안 2. 없는 API 경로 & 404 + Express 기본 HTML & 404 + JSON & \commit{38b4c12} \\
    \midrule
    N. 서버에서 도는 \code{clientRequest} & 키가 안 붙음, \code{ssr=false} & 키가 붙음, \code{ssr=true} & \commit{7f07df7} \\
    \midrule
    K. GCP 방화벽 & 22번, 3389번이 인터넷 전체에 열림 & 그대로 둠 (\decision{69}) & 없음 \\
    \bottomrule
  \end{tabularx}
\end{table}

%% ---------------------------------------------------------------------------
\section{CORS 거절을 403으로 (G)}\label{sec:05-cors}

\begin{inbrief}
  \item[Why] 아무 \texttt{Origin} 헤더로 DB 쓰기를 만들 수 있었고, 거절이 서버 고장(500)으로 기록됐다.
  \item[When] 허용 목록에 없는 \texttt{Origin}이 API를 부를 때다.
  \item[Where] \texttt{blog-api}의 Express 미들웨어(middleware)다.
  \item[How] 이름 있는 오류 \texttt{CorsRejectedError}를 던지고, \texttt{cors()} 바로 뒤에서 403으로 답한다.
\end{inbrief}

% src: deploy/traffic/review/03-status-codes.md §G (배포 완료 2026-10-13, 4d137f2)
\paragraph{결과.}
허용 안 된 \texttt{Origin}의 요청은 403과 WARN 한 줄로 끝난다.
DB 쓰기는 0건이다.
배포 뒤 Cloud Logging의 ERROR는 0건이었다.

\begin{figure}[H]
  \centering
  \begin{tikzpicture}[font=\small]
    \node[font=\small\bfseries] at (0,0.9) {전 (배포 전)};
    \node[font=\small\bfseries] at (7.2,0.9) {후 (\texttt{4d137f2})};
    \node[flowalt, minimum width=4.6cm] (b1) at (0,0) {요청\\(허용 안 된 Origin)};
    \node[flowalt, minimum width=4.6cm] (b2) at (0,-1.5) {CORS 거절};
    \node[flowalt, minimum width=4.6cm] (b3) at (0,-3.0) {errorHandler};
    \node[badbox, minimum width=4.6cm] (b4) at (0,-5.0) {ERROR 로그 2줄\\Supabase INSERT 1건\\500 응답};
    \draw[flowarrow] (b1) -- (b2);
    \draw[flowarrow] (b2) -- (b3);
    \draw[flowarrow] (b3) -- (b4);
    \node[flowbox, minimum width=4.6cm] (a1) at (7.2,0) {요청\\(허용 안 된 Origin)};
    \node[flowbox, minimum width=4.6cm] (a2) at (7.2,-1.5) {\texttt{CorsRejectedError}를 던짐};
    \node[flownew, minimum width=4.6cm] (a3) at (7.2,-3.0) {\texttt{corsRejection} 처리기\\(\texttt{cors()} 바로 뒤)};
    \node[goodbox, minimum width=4.6cm] (a4) at (7.2,-5.0) {403 응답, WARN 1줄\\DB 쓰기 없음};
    \draw[flowarrow] (a1) -- (a2);
    \draw[flowarrow] (a2) -- (a3);
    \draw[flowarrow] (a3) -- (a4);
    \draw[flowarrow, dashed] (a2.east) -- ++(0.9,0) |- (a4.east);
    \node[lbl, align=center, anchor=west] at (10.5,-3.2) {errorHandler\\안 거침};
  \end{tikzpicture}
  \flowlegend
  \caption{CORS 거절 한 건의 경로. 전은 errorHandler를 거쳐 500과 DB 쓰기가 생기고, 후는 처리기에서 끝난다.
  (src: deploy/traffic/review/03-status-codes.md §G)}\label{fig:05-cors-flow}
\end{figure}

\begin{description}
  \item[왜 필요했나] 주간 issue \#5에 ``blog-api ERROR 48'' 경고 표시가 떴다.
    원인을 찾다가 보안 문제를 발견했다.
  \item[문제] 거절이 \texttt{errorHandler}로 갔다.
    그래서 ERROR 로그 2줄, Supabase \texttt{error\_log} INSERT 1건, 500 응답이 생겼다.
    이것은 증폭 경로(amplification vector)다.
    아무 \texttt{Origin} 헤더로 요청을 많이 보내면 DB 쓰기가 같은 수만큼 생긴다.
    거절은 서버 고장이 아닌데 500으로 남아 주간 issue에 거짓 경고가 떴다.
    지난 7일 24건(ERROR 48줄)은 모두 우리가 만든 것이다. 로컬 설정이 23건, 고의 파손 시험이 1건이다.
  \item[무엇을 바꿨나] 파일 두 개다. \Cref{tab:05-cors-files}에 적었다.
  \item[어떻게] \texttt{server.ts}가 거절할 때 \texttt{CorsRejectedError}를 던진다.
    \texttt{cors()} 바로 뒤에 붙인 처리기가 이 오류만 잡아 403으로 답한다.
  \item[왜 이 방법인가] 아래 네 이유다.
\end{description}

\begin{itemize}[nosep]
  \item 이름 있는 오류를 쓴다. 처리기가 CORS 거절만 정확히 골라낸다. 진짜 오류는 그대로 \texttt{errorHandler}로 가서 기록된다.
  \item \texttt{cors()} 바로 뒤에 둔다. \texttt{errorHandler}보다 먼저 잡아야 DB 쓰기를 건너뛴다. 사전 요청(OPTIONS)도 여기서 잡힌다.
  \item 403이다. 요청한 쪽이 허락받지 않은 것이다. 서버 고장(5xx)이 아니다.
  \item 버린 대안: \texttt{callback(null, false)}로 조용히 넘기는 방법이다. 요청이 라우트까지 가서 DB를 읽는다. 브라우저만 막히고 서버 부하는 그대로다.
\end{itemize}

\begin{table}[H]
  \caption{G에서 바꾼 파일.}\label{tab:05-cors-files}
  % src: deploy/traffic/review/03-status-codes.md §G 무엇을 바꿨나
  \small
  \begin{tabularx}{\linewidth}{L{6.2cm} Y}
    \toprule
    파일 & 변경 \\
    \midrule
    \code{apps/dashboard_server/src/middlewares/corsRejection.ts} (새 파일) & \texttt{CorsRejectedError}. 이 오류만 잡아 403으로 답하는 처리기 \\
    \midrule
    \code{apps/dashboard_server/src/server.ts} & 거절할 때 \texttt{CorsRejectedError}를 던진다. \texttt{cors()} 바로 뒤에 처리기를 붙인다 \\
    \bottomrule
  \end{tabularx}
\end{table}

\begin{table}[H]
  \caption{G 전과 후, 요청 종류별.}\label{tab:05-cors-before-after}
  % src: deploy/traffic/review/03-status-codes.md §G 요청 표
  \small
  \begin{tabularx}{\linewidth}{Y Y Y}
    \toprule
    요청 & 전 & 후 \\
    \midrule
    허용 안 된 Origin, GET & 500, ERROR 2줄, INSERT 1건 & \textbf{403}, WARN 1줄, DB 쓰기 없음 \\
    \midrule
    허용 안 된 Origin, OPTIONS & 500 & \textbf{403} \\
    \midrule
    허용 Origin, Origin 없음 & 200 & 200 (그대로) \\
    \midrule
    다른 오류 & errorHandler & errorHandler (그대로) \\
    \bottomrule
  \end{tabularx}
\end{table}

\subsection{운영 확인 (2026/10/13 20:07 UTC)}\label{sec:05-cors-verify}

배포 전에는 타입 검사와 빌드가 통과했다.
배포 뒤 운영에서 \Cref{tab:05-cors-verify}의 여섯 가지를 확인했다.

\begin{table}[H]
  \caption{G 운영 확인. 대상 커밋 \texttt{4d137f2}.}\label{tab:05-cors-verify}
  % src: deploy/traffic/review/03-status-codes.md §G 확인 (운영, 2026-10-13 20:07 UTC)
  \small
  \begin{tabularx}{\linewidth}{L{6cm} Y}
    \toprule
    요청 & 결과 \\
    \midrule
    허용 안 된 Origin, OPTIONS & 403, \cfs{DYNAMIC} \\
    \midrule
    허용 안 된 Origin, GET \code{/users/me} & 403, \texttt{\{"success":false, "message":"origin not allowed"\}} \\
    \midrule
    허용 안 된 Origin, GET \code{posts} & 403, \cfs{EXPIRED} (캐시에서 그대로 나가지 않음) \\
    \midrule
    \texttt{Origin: my.}, OPTIONS / GET & 204 / 200 \cfs{HIT} \\
    \midrule
    Origin 없음 & 200 \cfs{HIT} \\
    \midrule
    Cloud Logging & 403은 WARNING \texttt{request completed} 한 줄씩. 배포 뒤 ERROR 0건 \\
    \bottomrule
  \end{tabularx}
\end{table}

확인에 쓴 명령은 한 줄이다. 기대값은 403이다.

\begin{lstlisting}[language=bash]
curl -s -o /dev/null -w '%{http_code}\n' -X OPTIONS -H 'Origin: https://evil.example' \
  -H 'Access-Control-Request-Method: GET' 'https://blog-api.prettylog.com/v1/api/posts?limit=1'   # 403
\end{lstlisting}

\begin{itemize}[nosep]
  \item GET은 Cloudflare 캐시가 먼저 200 \cfs{HIT}로 답할 수 있다. 서버가 요청을 받지 않으니 문제가 아니다.
  \item Cloud Logging에 \texttt{Unhandled Express error}가 새로 생기지 않는지 본다.
\end{itemize}

\subsection{CORS는 브라우저 규칙이다}\label{sec:05-cors-limit}

% src: deploy/traffic/review/02-blocking.md §오너 질문과 답 Q1
오너가 물었다: \texttt{ALLOWED\_ORIGINS}가 있는데 왜 키가 또 필요한가?
\code{curl} 같은 요청 때문인가?
답은 그렇다이다.
CORS(Cross-Origin Resource Sharing, 교차 출처 자원 공유)는 브라우저를 지키는 규칙이다.
서버를 지키는 자물쇠가 아니다.

\begin{table}[H]
  \caption{누구를 CORS가 막는가.}\label{tab:05-cors-who}
  % src: deploy/traffic/review/02-blocking.md §Q1
  \small
  \begin{tabularx}{\linewidth}{L{5cm} L{2.6cm} Y}
    \toprule
    요청한 쪽 & CORS가 막나 & 왜 \\
    \midrule
    다른 사이트에서 도는 브라우저 JS & 막는다 & 브라우저가 \texttt{Origin}을 붙이고, 응답 헤더를 보고 스스로 막는다 \\
    \midrule
    \code{curl}, 서버, 봇 & \textbf{못 막는다} & \texttt{Origin}을 안 붙이거나 마음대로 바꾼다. CORS를 지킬 의무가 없다 \\
    \bottomrule
  \end{tabularx}
\end{table}

\begin{itemize}[nosep]
  \item \texttt{Origin} 없는 요청을 통과시키는 이유: Next SSR, 헬스 체크, 같은 사이트 요청은 \texttt{Origin}을 보내지 않는다. 막으면 사이트가 깨진다.
  \item 그래서 DB에 쓰는 서버 전용 경로는 ``누가 보냈는지''를 증명하는 비밀 값으로 지킨다. 다음 절이 그 내용이다.
\end{itemize}

%% ---------------------------------------------------------------------------
\section{공개 error-logs 쓰기 막기 (L)}\label{sec:05-errorlogs}

\begin{inbrief}
  \item[Why] 인증 없는 \texttt{POST /v1/api/error-logs}가 요청 1건마다 Supabase에 1행을 썼다.
  \item[When] 인터넷의 누구나 \code{curl}로 부를 수 있을 때다.
  \item[Where] \texttt{blog-api}의 error-logs 라우트와 Next 서버의 호출 세 곳이다.
  \item[How] Next 서버가 \texttt{x-origin-key} 비밀 헤더를 보내고, Express가 비교한다. 틀리면 403이다.
\end{inbrief}

% src: deploy/traffic/review/03-status-codes.md §L (배포 완료 2026-10-13 38b4c12, 결정 #60)
\paragraph{결과.}
키 없이 보낸 POST는 403 \texttt{forbidden}이다.
틀린 키도 403이다.
배포 뒤 ERROR는 0건이었다.

\begin{figure}[H]
  \centering
  \begin{tikzpicture}[font=\small]
    \node[flowbox, minimum width=2.6cm] (br) at (0,0) {브라우저 오류};
    \node[flowbox, minimum width=2.8cm] (nx) at (3.7,0) {Next\\\texttt{/api/log-error}};
    \node[flownew, minimum width=3.4cm] (ex) at (7.9,0) {Express \texttt{/error-logs}\\\texttt{requireOriginKey}};
    \node[dbbox, minimum width=2cm] (db) at (11.9,0) {INSERT};
    \draw[flowarrow] (br) -- (nx);
    \draw[flowarrow] (nx) -- (ex) node[lbl, midway, above] {\texttt{x-origin-key}};
    \draw[flowarrow] (ex) -- (db) node[lbl, midway, above] {키 일치};
    \node[flowalt, minimum width=2.6cm] (cu) at (0,-2.2) {누구나 \texttt{curl}};
    \node[badbox, minimum width=2.4cm] (no) at (11.9,-2.2) {403, INSERT 없음};
    \draw[flowarrow] (cu) -- (ex.south west |- cu) node[lbl, midway, above] {헤더 없음} -- (ex.south west);
    \draw[flowarrow] (ex.south east) |- (no) node[lbl, near end, above] {키 없음 또는 틀림};
  \end{tikzpicture}
  \flowlegend
  \caption{error-logs 쓰기 경로. 위는 Next 서버가 키를 붙여 통과하고, 아래는 키 없는 \texttt{curl}이 403으로 끝난다.
  (src: deploy/traffic/review/03-status-codes.md §L)}\label{fig:05-originkey-gate}
\end{figure}

\begin{description}
  \item[왜 필요했나] 설계 검토 중 architect가 ``인증 여부 확인 필요''로 남긴 경로를 확인했다.
  \item[문제] \texttt{POST /v1/api/error-logs}에는 인증이 없다.
    인터넷 누구나 부를 수 있다.
    요청 1건마다 Supabase에 1행을 쓴다.
    CORS 문제(G)보다 직접적인 DB 쓰기 경로다.
    CORS는 브라우저만 막고 \code{curl}은 막지 못한다.
  \item[무엇을 바꿨나] 파일 네 묶음이다. \Cref{tab:05-originkey-files}에 적었다.
  \item[어떻게] VM의 \texttt{shared/api.env}와 \texttt{shared/web.env}에 같은 \texttt{ORIGIN\_KEY}를 넣는다.
    64자이고 \texttt{openssl rand -hex 32}로 만든다.
    웹을 먼저, API를 나중에 올린다.
    둘 다 ecosystem 파일로 restart한다.
  \item[왜 이 방법인가] 아래 다섯 이유다.
\end{description}

\begin{itemize}[nosep]
  \item 공유 비밀이다. 이 경로는 Next 서버만 부른다. 사용자 로그인과 상관없다. 서버끼리의 증명에는 비밀 헤더 하나면 된다.
  \item 비교에 \texttt{timingSafeEqual}을 쓴다. 응답 시간 차이로 키를 한 글자씩 알아내는 공격을 막는다.
  \item 키가 없으면 예전처럼 연다. 배포 순서가 틀려도 사이트와 오류 기록이 멈추지 않는다. 시작할 때 \texttt{ORIGIN\_KEY unset} 경고를 남겨 빠뜨린 것을 알 수 있다.
  \item 요청 수 제한만으로 막지 않는다. 한도 안의 스팸은 그대로 쌓인다 (\decision{60}).
  \item 같은 헤더를 SSR 제외(D12, \decision{64})에 다시 쓴다. 비밀을 하나만 관리한다.
\end{itemize}

\begin{callout}[주의]
\texttt{ORIGIN\_KEY=}처럼 빈 값을 두면 Express가 시작하지 못한다(최소 32자 검사).
쓰지 않을 때는 줄을 지운다.
\end{callout}

\begin{table}[H]
  \caption{L에서 바꾼 파일.}\label{tab:05-originkey-files}
  % src: deploy/traffic/review/03-status-codes.md §L 무엇을 바꿨나
  \small
  \begin{tabularx}{\linewidth}{L{6.4cm} Y}
    \toprule
    파일 & 변경 \\
    \midrule
    \code{apps/dashboard_server/src/middlewares/requireOriginKey.ts} (새 파일) & \texttt{x-origin-key}가 env \texttt{ORIGIN\_KEY}와 같을 때만 통과. 아니면 403 \\
    \midrule
    \code{apps/dashboard_server/src/controllers/index.ts} & \texttt{/api/error-logs} 앞에 붙임 \\
    \midrule
    \code{apps/dashboard_client_web/apis/http.server.ts} & \texttt{originKeyHeaders()}. 키가 있을 때만 헤더를 만든다 \\
    \midrule
    \code{instrumentation.ts}, \code{app/api/log-error/route.ts}, \code{apis/reportSsrError.server.ts} & \texttt{/error-logs}를 부르는 세 곳만 헤더를 보낸다 \\
    \bottomrule
  \end{tabularx}
\end{table}

\subsection{운영 확인 (2026/10/13 21:0x UTC)}\label{sec:05-errorlogs-verify}

% src: deploy/traffic/review/03-status-codes.md §L 확인 (2026-10-13 21:0x UTC, 38b4c12)
\Cref{tab:05-originkey-verify}에 대상 커밋 \commit{38b4c12}의 확인 결과를 적었다.

\begin{table}[H]
  \caption{L 운영 확인.}\label{tab:05-originkey-verify}
  % src: deploy/traffic/review/03-status-codes.md §L 확인
  \small
  \begin{tabularx}{\linewidth}{L{6.4cm} Y}
    \toprule
    확인 & 결과 \\
    \midrule
    키 없이 POST & 403 \texttt{forbidden} \\
    \midrule
    틀린 키 & 403 \\
    \midrule
    \code{/v1/api/nope} & 404 JSON \\
    \midrule
    두 호스트의 \texttt{x-powered-by} & 없음 \\
    \midrule
    \code{/}, \code{/blog}, 글, \code{/about} & 200 \\
    \midrule
    \code{traffic-check.sh} & 0 FAIL \\
    \midrule
    배포 뒤 ERROR & 0건 \\
    \midrule
    \texttt{ORIGIN\_KEY unset} 경고 & 21:01:13에 1번 (\code{release.sh}의 reload가 키를 읽기 전). 그 뒤 restart부터는 없음 \\
    \bottomrule
  \end{tabularx}
\end{table}

%% ---------------------------------------------------------------------------
\section{서버에서 도는 clientRequest도 키를 보낸다 (N)}\label{sec:05-clientrequest}

\begin{inbrief}
  \item[Why] Express 요청 수 제한(L-b)은 \texttt{ssr=true} 요청을 세지 않는다. 모든 SSR 요청이 \texttt{ssr=true}여야 한다.
  \item[When] \texttt{enforce}로 바꾸기 전이다.
  \item[Where] \code{apis/http.client.ts}의 \texttt{clientRequest}다.
  \item[How] 서버에서 도는지(\texttt{typeof window})를 보고, 서버일 때만 \texttt{x-origin-key}를 붙인다.
\end{inbrief}

% src: deploy/traffic/review/02-blocking.md §N (웹 배포 완료 7f07df7, 결정 #64)
\paragraph{결과.}
MISS 페이지 5개를 연 뒤 blog-api 로그의 \texttt{categories} 8건, \texttt{comments} 1건, \texttt{/v1/api/log} 1건이 모두 \texttt{ssr=true}였다.

\begin{description}
  \item[왜 필요했나] Next SSR 요청은 모두 VM IP 하나로 보인다.
    키가 없으면 IP로 세어져 사이트가 자기 SSR을 막는다.
  \item[문제] 배포 확인 중 \texttt{categories} SSR 호출이 \texttt{ssr=false}로 찍혔다.
    일부 fetcher는 서버와 브라우저 양쪽에서 쓴다(isomorphic).
    이 fetcher들은 브라우저용 \texttt{clientRequest}를 쓰고, 서버에서 불려도 키가 붙지 않았다.
    해당 fetcher는 \Cref{tab:05-fetchers}의 다섯 개다.
  \item[무엇을 바꿨나] \code{apis/http.client.ts}의 \texttt{clientRequest}가 지금 서버에서 도는지를 보고 헤더를 정한다.
  \item[어떻게] \Cref{fig:05-clientrequest} 아래 네 걸음이다.
\end{description}

\begin{enumerate}[nosep]
  \item 같은 파일이 두 번 실행된다. 한 번은 Next 서버(Node)에서, 한 번은 사용자 브라우저에서다.
  \item Node에는 \texttt{window} 객체가 없다. 브라우저에는 있다. 그래서 \texttt{typeof window}로 어디서 도는지 안다.
  \item 서버에서는 \texttt{process.env.ORIGIN\_KEY}를 읽어 헤더를 붙인다. 값은 VM의 \texttt{shared/web.env}에서 온다.
  \item Express는 그 헤더를 \texttt{ORIGIN\_KEY}와 비교한다. 같으면 \texttt{ssr=true}로 기록하고 요청 수를 세지 않는다.
\end{enumerate}

\begin{figure}[H]
  \centering
  \begin{tikzpicture}[font=\small]
    \node[flownew, minimum width=3.4cm] (cr) at (0,0) {\texttt{clientRequest} 호출};
    \node[flowbox, minimum width=3.6cm] (br) at (-3.6,-2.0) {\texttt{window} 있음\\(브라우저)};
    \node[flowbox, minimum width=3.6cm] (sv) at (3.6,-2.0) {\texttt{window} 없음\\(Next 서버)};
    \node[flowbox, minimum width=3.6cm] (b2) at (-3.6,-4.0) {헤더 없음\\사용자 IP로 세어짐};
    \node[goodbox, minimum width=3.6cm] (s2) at (3.6,-4.0) {\texttt{ORIGIN\_KEY}가 있으면\\\texttt{x-origin-key} 붙임\\\texttt{ssr=true}};
    \draw[flowarrow] (cr) -- (br);
    \draw[flowarrow] (cr) -- (sv);
    \draw[flowarrow] (br) -- (b2);
    \draw[flowarrow] (sv) -- (s2);
  \end{tikzpicture}
  \caption{\texttt{clientRequest}가 실행 장소에 따라 헤더를 정하는 길.
  브라우저에는 헤더가 붙지 않는다. (src: deploy/traffic/review/02-blocking.md §N)}\label{fig:05-clientrequest}
\end{figure}

\begin{table}[H]
  \caption{키를 안 보내던 isomorphic fetcher와 서버에서 부르는 곳.}\label{tab:05-fetchers}
  % src: deploy/traffic/review/02-blocking.md §N fetcher 표
  \small
  \begin{tabularx}{\linewidth}{L{4.4cm} Y}
    \toprule
    fetcher & 서버에서 부르는 곳 \\
    \midrule
    \code{getCategories.api.ts} & \code{CategorySidebarWrapper.server.tsx}, \code{sidebar-categories.server.tsx}, \code{getFilterItems.server.ts} \\
    \midrule
    \code{getComments.api.ts} & \code{Comments.server.tsx} \\
    \midrule
    \code{get-log-thread.api.ts} & \code{app/log/[username]/[logId]/page.tsx} \\
    \midrule
    \code{get-log-global.api.ts} & \code{app/log/page.tsx} \\
    \midrule
    \code{get-log-stream.api.ts} & \code{app/log/[username]/page.tsx} \\
    \bottomrule
  \end{tabularx}
\end{table}

\paragraph{왜 이 방법인가.}
\begin{itemize}[nosep]
  \item \code{http.server.ts}의 함수를 다시 쓰지 않았다. 그 파일은 \texttt{import 'server-only'}가 있다. 브라우저에도 들어가는 \code{http.client.ts}가 import하면 빌드가 깨진다. 그래서 같은 4줄을 \code{http.client.ts}에 두었다.
  \item fetcher 5개를 따로 고치지 않았다. \texttt{clientRequest} 한 곳만 고치면 지금 것과 앞으로 생길 isomorphic fetcher가 모두 덮인다.
  \item 브라우저 요청에는 영향이 없다. 헤더가 안 붙으므로 CORS 사전 요청(preflight)도 늘지 않는다.
\end{itemize}

\paragraph{왜 브라우저로 키가 새지 않나.}
\Cref{tab:05-leak-layers}의 두 겹이다.

\begin{table}[H]
  \caption{브라우저로 키가 새지 않는 두 겹.}\label{tab:05-leak-layers}
  % src: deploy/traffic/review/02-blocking.md §N 왜 브라우저로 키가 새지 않나
  \small
  \begin{tabularx}{\linewidth}{l Y}
    \toprule
    겹 & 이유 \\
    \midrule
    1. 이름 & \texttt{ORIGIN\_KEY}는 \texttt{NEXT\_PUBLIC\_}으로 시작하지 않는다. Next는 이런 변수를 브라우저 코드에 넣지 않는다. 브라우저에서는 값이 \texttt{undefined}다 \\
    \midrule
    2. 조건 & 브라우저에서는 \texttt{window}가 있으므로 키를 읽는 줄까지 가지 않는다 \\
    \bottomrule
  \end{tabularx}
\end{table}

\subsection{확인}\label{sec:05-clientrequest-verify}

\begin{itemize}[nosep]
  \item 2026/10/13 21:20 UTC, 웹 \commit{7f07df7}. MISS 페이지 5개를 열었다 (\code{/blog/@taeklim}, 글 하나, \code{/log}).
  \item 결과: \texttt{categories} 8건, \texttt{comments} 1건, \texttt{/v1/api/log} 1건이 모두 \texttt{ssr=true}였다.
  \item 확인하지 못한 것: 빌드된 브라우저 번들에 키 문자열이 없는지 직접 여는 확인이다. 로컬 \texttt{next build}가 금지라서 하지 않았다. 배포 뒤 개발자 도구에서 API 요청 헤더에 \texttt{x-origin-key}가 없는지 보면 된다. 이 확인은 오너 몫이다.
  \item 아래 쿼리는 브라우저 호출만 남겨야 한다. \texttt{ssr=true}로 바꾸면 SSR 호출이 보여야 한다.
\end{itemize}

\begin{lstlisting}
log_id("blog-api") AND jsonPayload.req.url:"/categories" AND jsonPayload.req.ssr=false
\end{lstlisting}

%% ---------------------------------------------------------------------------
\section{서버 종류 헤더 숨기기와 JSON 404}\label{sec:05-headers}

\begin{inbrief}
  \item[Why] 공격자에게 프레임워크를 알려 줄 이유가 없고, 응답 형식은 하나여야 한다.
  \item[When] 2026/10/13 측정에서 \texttt{x-powered-by}와 Express 기본 404 HTML이 보였을 때다.
  \item[Where] \texttt{next.config.ts}와 Express의 마지막 404 처리기다.
  \item[How] 헤더를 끄고, 없는 API 경로는 JSON으로 답한다.
\end{inbrief}

% src: deploy/traffic/review/03-status-codes.md §제안 1, 2 (배포 완료 38b4c12)
\paragraph{결과.}
두 호스트에서 \texttt{x-powered-by}가 없어졌다.
\code{/v1/api/nope}는 404 JSON \texttt{\{success:false\}}를 돌려준다.
\Cref{tab:05-proposals}의 1, 2번이 이 내용이다.

\begin{table}[H]
  \caption{오늘 측정으로 찾은 제안 네 개. 1, 2는 코드 완료, 3, 4는 설계가 먼저다.}\label{tab:05-proposals}
  % src: deploy/traffic/review/03-status-codes.md §제안
  \scriptsize\setlength{\tabcolsep}{2.5pt}
  \begin{tabularx}{\linewidth}{L{1.7cm} Y Y Y Y}
    \toprule
    \# & 지금 (측정 10/13) & 바꾸면 & 왜 & 어디서 \\
    \midrule
    1 (배포 완료 \commit{38b4c12}) & \texttt{x-powered-by:} \texttt{Next.js}, \texttt{x-powered-by:} \texttt{Express} & 헤더 없음 & 공격자에게 프레임워크를 알려 줄 이유가 없다 & \texttt{next.config.ts} \texttt{poweredByHeader:} \texttt{false}, Express \code{app.disable('x-powered-by')} \\
    \midrule
    2 (배포 완료 \commit{38b4c12}) & 404 + Express 기본 HTML (\texttt{Cannot GET /v1/api/nope}) & 404 + JSON \texttt{\{success:false\}} & 응답 형식을 하나로. 내부 경로 모양을 덜 드러낸다 & Express 마지막 404 처리기 \\
    \midrule
    3 보안 헤더 (미정) & HSTS, \code{X-Content-Type-Options}, \code{Referrer-Policy}, \code{X-Frame-Options} 없음 (API 404에만 일부) & 모든 응답에 붙음 & 다운그레이드, MIME 추측, 클릭재킹(clickjacking) 방지 & Cloudflare Transform Rules 또는 앱. HSTS는 Cloudflare SSL/TLS 설정 (preload는 되돌리기 어려우니 빼고 시작) \\
    \midrule
    4 피드 실패 화면 (미정) & 200, 캐시 가능 & 그 응답만 짧게 또는 캐시 안 함 & 일시 장애가 하루 동안 남지 않게 & 설계 필요 (Next에서 응답 헤더를 바꾸기 어렵다) \\
    \bottomrule
  \end{tabularx}
\end{table}

%% ---------------------------------------------------------------------------
\section{상태 코드 표}\label{sec:05-status-table}

\begin{inbrief}
  \item[Why] 응답 코드가 의미를 말해야 캐시, 로그, 경고가 맞게 동작한다.
  \item[When] 새 경로나 오류 처리를 추가할 때 이 표를 기준으로 삼는다.
  \item[Where] Next 페이지, Express API, Cloudflare 오류 응답이다.
  \item[How] 거절은 4xx, 서버 고장만 5xx로 하고, 오류 응답은 캐시하지 않는다.
\end{inbrief}

% src: deploy/traffic/review/03-status-codes.md §이미 있는 상태 코드 정리
\Cref{tab:05-status-codes}는 이미 정해진 경우다.

\begin{table}[H]
  \caption{경우별 상태 코드와 캐시 여부.}\label{tab:05-status-codes}
  % src: deploy/traffic/review/03-status-codes.md §이미 있는 상태 코드 정리
  \small
  \begin{tabularx}{\linewidth}{L{4cm} L{2.8cm} L{2.4cm} Y}
    \toprule
    경우 & 코드 & 캐시 & 왜 \\
    \midrule
    없는 글 (\code{getPost.api.ts}, \commit{42ba55c}) & 404 페이지 & 안 함 (300--599 규칙) & 예전에는 오류 화면(500 경계)이 떴다. 없는 것은 404가 맞다 \\
    \midrule
    피드 API가 죽었을 때 (\commit{77315eb}) & 200 + ``잠시 뒤 다시'' 안내 & 1일 캐시 가능 (경고) & 화면 전체가 오류로 바뀌지 않게 했다. 아래 주의 참고 \\
    \midrule
    보호 경로, 로그인 안 함 (\code{/me}) & 307 $\rightarrow$ \code{/sign-in} & 안 함 (bypass) & Clerk \texttt{auth.protect} \\
    \midrule
    API 로그인 필요, 토큰 없음 & 401 & 안 함 (\cfs{EXPIRED}) &  \\
    \midrule
    Cloudflare 오류 응답 & 300--599 & 안 함 & 서버가 잠깐 죽은 순간의 화면이 하루 동안 남지 않게 \\
    \bottomrule
  \end{tabularx}
\end{table}

\begin{callout}[주의: 피드 안내 화면]
피드를 못 받은 홈은 200으로 나간다.
그래서 ``web cache'' 규칙이 하루 동안 캐시할 수 있다.
이때는 purge 한 번으로 고친다.
근본 해결은 \Cref{tab:05-proposals}의 제안 4다.
\end{callout}

%% ---------------------------------------------------------------------------
\section{GCP 방화벽 (K, 중단)}\label{sec:05-firewall}

\begin{inbrief}
  \item[Why] Express backstop은 헤더가 없으면 제한하지 않는다(fail-open). 이것은 3000, 4000번 포트가 인터넷에 닫혀 있을 때만 안전하다.
  \item[When] 2026/10/13에 \texttt{gcloud compute firewall-rules list}로 확인했다.
  \item[Where] GCP 프로젝트의 방화벽(firewall) 규칙이다.
  \item[How] 기본 규칙 둘을 끄는 안이 있었지만 오너가 중단했다 (\decision{69}).
\end{inbrief}

% src: deploy/traffic/review/02-blocking.md §K (중단, 오너 #69)
\paragraph{결과.}
fail-open의 전제는 맞다.
3000, 4000번은 열려 있지 않다.
그러나 기본 규칙 두 개가 22번과 3389번을 인터넷 전체에 열어 둔다.
오너가 ``it is ok. gcp already taken care of''라고 답해 그대로 두기로 했다.

\begin{table}[H]
  \caption{방화벽 규칙 확인 (2026/10/13).}\label{tab:05-firewall}
  % src: deploy/traffic/review/02-blocking.md §K 확인
  \small
  \begin{tabularx}{\linewidth}{L{4cm} Y L{3.6cm}}
    \toprule
    규칙 & 열린 것 & 판정 \\
    \midrule
    3000, 4000번 & 규칙 없음. 앱도 \texttt{127.0.0.1}에서만 받음 & 안전. fail-open 전제가 맞다 \\
    \midrule
    \code{allow-ssh-from-iap} & 22번, IAP 대역 \texttt{35.235.240.0/20}만 & 필요. 우리 접속 경로다 \\
    \midrule
    \code{default-allow-ssh} & 22번, \textbf{인터넷 전체} & 필요 없음 \\
    \midrule
    \code{default-allow-rdp} & 3389번(Windows 원격), \textbf{인터넷 전체} & 필요 없음. VM은 Linux다 \\
    \bottomrule
  \end{tabularx}
\end{table}

\begin{description}
  \item[무엇을 바꾸나] 두 기본 규칙을 끈다. 지우지 않고 끄기만 해서 되돌릴 수 있다.
  \item[왜 이 방법인가] 공개 22번은 비밀번호 대입과 취약점 탐색의 첫 대상이다. 우리는 IAP로만 들어가므로 잃는 것이 없다.
  \item[상태] 중단이다. 명령은 \Cref{ch:commands}의 방화벽 절에 있다. 오너가 켜기로 하면 그때 쓴다.
\end{description}

%% ---------------------------------------------------------------------------
\section{비밀 값과 개인 정보}\label{sec:05-secrets}

\begin{inbrief}
  \item[Why] 저장소가 공개라서 비밀 값과 개인 정보를 한 줄도 넣을 수 없다.
  \item[When] 문서, 스크립트, issue를 쓸 때마다다.
  \item[Where] 저장소, GitHub issue, 터미널 기록이다.
  \item[How] \Cref{tab:05-secrets}의 규칙을 따른다.
\end{inbrief}

% src: deploy/traffic/review/05-secrets.md §비밀 값과 개인 정보 다루기
\begin{table}[H]
  \caption{비밀 값과 개인 정보를 다루는 규칙.}\label{tab:05-secrets}
  % src: deploy/traffic/review/05-secrets.md §비밀 값과 개인 정보 다루기
  \small
  \begin{tabularx}{\linewidth}{L{3.8cm} Y Y}
    \toprule
    항목 & 어떻게 다루나 & 왜 \\
    \midrule
    방문자 IP & Cloud Logging에만 저장. 30일 뒤 지워짐 & 공격 대응에 필요. 저장소는 공개다 \\
    \midrule
    결제 계정 ID, 메일 주소, 토큰 & 저장소와 issue에 넣지 않음. 문서에는 \code{<ACCOUNT_ID>} & 공개 저장소 \\
    \midrule
    메일 주소 입력 & \code{read}로 입력. 터미널 변수에만 있음 & 셸 기록과 파일에 남지 않게 \\
    \midrule
    Cloudflare purge 토큰 & 환경 변수 또는 \code{_scripts/cloudflare/.env} (\code{*.env}는 gitignore) & 커밋되지 않게 \\
    \midrule
    주간 issue & 캐시 검사 줄만 넣음 & 공개 URL과 헤더뿐 \\
    \bottomrule
  \end{tabularx}
\end{table}

\begin{itemize}[nosep]
  \item 이 책의 부록도 같은 규칙을 따른다. 계정 ID, IP, 메일은 \texttt{<PLACEHOLDER>} 형태로만 적었다.
  \item \texttt{ORIGIN\_KEY}는 VM의 \texttt{shared/api.env}, \texttt{shared/web.env}에만 있다. 저장소, 브라우저, 로그에는 없다.
  \item 키가 새면 VM에서 새 키를 만들어 두 파일을 바꾸고 두 앱을 restart한다.
\end{itemize}
